\subsection{INVERSE LIMITS OF MODULES}

Let \((\mathbf{A}_{\alpha}, \phi_{\alpha\beta})\) be an inverse system of rings (I, § 10, no. 1), \((\mathbf{E}_{\alpha}, f_{\alpha\beta})\) an inverse system of commutative groups (written additively) (I, § 10, no. 1) and suppose that each \(E_{\alpha}\) has a left \(A_{\alpha}\) -module structure; moreover suppose that for \(\alpha \leqslant \beta\) \((f_{\alpha\beta}, \phi_{\alpha\beta})\) is a dimorphism of \(E_{\beta}\) into \(E_{\alpha}\) (§ 1, no. 13), in other words that

\[
f _ {\alpha \beta} (\lambda_ {\beta} x _ {\beta}) = \phi_ {\alpha \beta} (\lambda_ {\beta}) f _ {\alpha \beta} (x _ {\beta}),\tag{1}
\]

for \(x_{\beta} \in \mathrm{E}_{\beta}, \lambda_{\beta} \in \mathrm{A}_{\beta}\) ; then it follows from I, § 10, no. 2 that \(\mathbf{E} = \varprojlim \mathbf{E}_{\alpha}\) has a left module structure over \(\mathbf{A} = \varprojlim \mathbf{A}_{\alpha}\) . For all \(\alpha \in \mathbf{I}\) , let \(f: \mathbf{E} \to \mathbf{E}_{\alpha}, \phi_{\alpha}: \mathbf{A} \to \mathbf{A}_{\alpha}\) be the canonical mappings; then \((f_{\alpha}, \phi_{\alpha})\) is a dimorphism of \(\mathbf{E}\) into \(\mathbf{E}_{\alpha}\) . We shall say that \((\mathbf{E}_{\alpha}, f_{\alpha\beta})\) is an inverse system of left \(A_{\alpha}\) -modules and that the A-module E is its inverse limit.

Let \((\mathbf{E}_{\alpha}^{\prime}, f_{\alpha\beta}^{\prime})\) be another inverse system of left \(A_{\alpha}\) -modules and, for all \(\alpha\) , let \(u_{\alpha}: E_{\alpha}^{\prime} \to E_{\alpha}\) be an \(A_{\alpha}\) -linear mapping, these mappings forming an inverse system; then \(u = \lim_{\leftarrow} u_{\alpha}\) is an A-linear mapping of \(\lim_{\leftarrow} E_{\alpha}^{\prime}\) into \(\lim_{\leftarrow} E_{\alpha}\) .

Moreover:

PROPOSITION 1. Let \((\mathbf{E}_{\alpha}, f_{\alpha\beta})\) , \((\mathbf{E}_{\alpha}^{\prime}, f_{\alpha\beta}^{\prime})\) , \((\mathbf{E}_{\alpha}^{\prime \prime}, f_{\alpha\beta}^{\prime \prime})\) be three inverse systems of \(\mathbf{A}_{\alpha}\) -modules and \((u_{\alpha})\) , \((v_{\alpha})\) two inverse systems of \(\mathbf{A}_{\alpha}\) -linear mappings such that the sequences

\[
0 \longrightarrow \mathrm{E} _ {\alpha} ^ {\prime} \stackrel {u _ {\alpha}} {\longrightarrow} \mathrm{E} _ {\alpha} \stackrel {v _ {\alpha}} {\longrightarrow} \mathrm{E} _ {\alpha} ^ {\prime \prime}
\]

are exact for all \(\alpha\) . Then, writing \(u = \varprojlim u_{\alpha}\) , \(v = \varprojlim v_{\alpha}\) , the sequence

\[
0 \longrightarrow \varprojlim \mathrm{E} _ {\alpha} ^ {\prime} \stackrel {u} {\longrightarrow} \varprojlim \mathrm{E} _ {\alpha} \stackrel {v} {\longrightarrow} \varprojlim \mathrm{E} _ {\alpha} ^ {\prime \prime}
\]

is exact.

As \(u_{\alpha}^{-1}(0) = \{0\}\) for all \(\alpha\) , it follows from Set Theory, III, § 7, no. 2, Proposition 2 that \(u^{-1}(0) = \{0\}\) , hence \(u\) is injective; further, the \(u_{\alpha}(\mathrm{E}_{\alpha}')\) form an inverse system of subsets of the \(\mathrm{E}_{\alpha}\) and thus \(u(\varprojlim \mathrm{E}_{\alpha}') = \varprojlim u_{\alpha}(\mathrm{E}_{\alpha}')\) . As \(u_{\alpha}(\mathrm{E}_{\alpha}') = v_{\alpha}^{-1}(0)\) by hypothesis, \(v^{-1}(0) = \varprojlim u_{\alpha}(\mathrm{E}_{\alpha}') = u(\varprojlim \mathrm{E}_{\alpha}')\) (Set Theory, III, § 7, no. 2, Proposition 2), which completes the proof.

Remarks. (1) Proposition 1 and its proof are valid for arbitrary groups, except for change of notation.

\begin{enumerate}
\def\labelenumi{(\arabic{enumi})}
\setcounter{enumi}{1}
\tightlist
\item
  Note that if there are exact sequences
\end{enumerate}

\[
0 \longrightarrow \mathrm{E} _ {\alpha} ^ {\prime} \stackrel {u _ {\alpha}} {\longrightarrow} \mathrm{E} _ {\alpha} \stackrel {v _ {\alpha}} {\longrightarrow} \mathrm{E} _ {\alpha} ^ {\prime \prime} \longrightarrow 0
\]

it does not necessarily follow that the sequence

\[
0 \longrightarrow \varprojlim \mathrm{E} _ {\alpha} ^ {\prime} \stackrel {u} {\longrightarrow} \varprojlim \mathrm{E} _ {\alpha} \stackrel {v} {\longrightarrow} \varprojlim \mathrm{E} _ {\alpha} ^ {\prime \prime} \longrightarrow 0
\]

is exact; in other words, the inverse limit of an inverse system of surjective linear mappings is not necessarily surjective (cf.~Exercise 1).

Suppose now that the \(\mathbf{A}_{\alpha}\) are equal to the same ring \(\mathbf{A}\) and the \(\phi_{\alpha \beta}\) to \(1_{\mathbf{A}}\) ; then for every inverse system \((\mathrm{E}_{\alpha}, f_{\alpha \beta})\) of A-modules, \(\mathbf{E} = \varprojlim \mathbf{E}_{\alpha}\) is an A-module. Let \(\mathbf{F}\) be an A-module and, for all \(\alpha\) , let \(u_{\alpha}: \mathbf{F} \to \mathbf{E}_{\alpha}\) be an A-linear mapping such that \((u_{\alpha})\) is an inverse system of mappings; then \(u = \varprojlim u_{\alpha}\) is an A-linear mapping of \(\mathbf{F}\) into \(\varprojlim \mathbf{E}_{\alpha}\) . Conversely, for every A-linear mapping \(v: \mathbf{F} \to \varprojlim \mathbf{E}_{\alpha}\) , the family of \(v_{\alpha} = f_{\alpha} \circ v\) is an inverse system of A-linear mappings such that \(v = \varprojlim v_{\alpha}\) . We note on the other hand that for \(\alpha \leqslant \beta\) the mapping

\[
\operatorname{Hom} \left(\mathrm{l} _ {\mathrm{F}}, f _ {\alpha \beta}\right) = \bar {f} _ {\alpha \beta}: \operatorname{Hom} _ {\mathrm{A}} (\mathrm{F}, \mathrm{E} _ {\beta}) \rightarrow \operatorname{Hom} _ {\mathrm{A}} (\mathrm{F}, \mathrm{E} _ {\alpha})
\]

is a Z-module homomorphism such that \((\mathrm{Hom}_{\mathbf{A}}(\mathrm{F}, \mathrm{E}_{\alpha}), \vec{f}_{\alpha\beta})\) is an inverse system of Z-modules; as \(\vec{f}_{\alpha\beta}(v_{\beta}) = f_{\alpha\beta} \circ v_{\beta}\) , the above remarks can therefore be expressed as follows:

PROPOSITION 2. For every inverse system \((\mathbf{E}_{\alpha}, f_{\alpha \beta})\) of A-modules and every A-module F, the canonical mapping \(u \mapsto (f_{\alpha} \circ u)\) is a Z-module isomorphism

\[
l _ {\mathrm{F}}: \operatorname{Hom} _ {\mathrm{A}} (\mathrm{F}, \varprojlim \mathrm{E} _ {\alpha}) \rightarrow \varprojlim \operatorname{Hom} _ {\mathrm{A}} (\mathrm{F}, \mathrm{E} _ {\alpha}).\tag{2}
\]

COROLLARY. For every A-module homomorphism \(v: \mathbf{F} \to \mathbf{F}'\) , the

\[
\bar {v} _ {\alpha} = \operatorname{Hom} (v, 1 _ {\mathrm{E} _ {\alpha}}): \operatorname{Hom} (\mathrm{F} ^ {\prime}, \mathrm{E} _ {\alpha}) \rightarrow \operatorname{Hom} (\mathrm{F}, \mathrm{E} _ {\alpha})
\]

form an inverse system of Z-linear mappings and the diagram

\[
\begin{array}{c} \operatorname{Hom} (\mathrm{F} ^ {\prime}, \varprojlim \mathrm{E} _ {\alpha}) \xrightarrow {l _ {\mathrm{F} ^ {\prime}}} \varprojlim \operatorname{Hom} (\mathrm{F} ^ {\prime}, \mathrm{E} _ {\alpha}) \\ \operatorname{Hom} (v, 1 _ {\mathrm{E}}) \Biggl \downarrow \qquad \qquad \qquad \qquad \qquad \qquad \qquad \qquad \qquad \qquad \qquad \qquad \qquad \qquad \qquad \qquad \qquad \qquad \qquad \qquad \qquad \qquad \qquad \qquad \qquad \qquad \qquad \qquad \qquad \qquad \qquad \qquad \qquad \qquad \\ \operatorname{Hom} (\mathrm{F}, \varprojlim \mathrm{E} _ {\alpha}) \xrightarrow {l _ {\mathrm{F}}} \varprojlim \operatorname{Hom} (\mathrm{F}, \mathrm{E} _ {\alpha}) \end{array}\tag{3}
\]

is commutative.

For all \(u \in \operatorname{Hom}(\mathbf{F}', \varprojlim \mathbf{E}_{\alpha})\) , \(l_{\mathbf{F}}(u \circ v) = (f_{\alpha} \circ u \circ v)\) by definition and the commutativity of diagram (3) follows immediately from the definitions.

